Le rétroviseur extérieur d'une voiture est un miroir convexe dont le rayon de courbure est de \qty{\v{R}}{m}. Une auto de \qty{\v{ho}}{m} de hauteur suit la voiture et se trouve à \qty{\v{dob}}{m} du rétroviseur.

\begin{enumerate}[label=\alph*)]
    \item Déterminez la distance focale du rétroviseur (attention au signe).
    \item Où se forme l'image de l'auto ?
    \item Déterminez la hauteur de l'image et décrivez ses caractéristiques. Pourquoi ces rétroviseurs portent-ils la mention « Les objets sont plus proches qu'ils ne le paraissent » ?
\end{enumerate}
